% !TeX root = ../初等数论笔记.tex
% 由主文档通过 \input 载入；请编译主文档。
\setcounter{chapter}{13} % 保留原书章号。
\chapter{数论算法与公钥密码入门（选学）}
\label{ch:algorithms}

前面的定理不仅回答“为什么”，也提供了计算步骤。把定理变成算法，还要规定输入范围、输出含义、何时终止，以及计算量怎样随输入增长。本章以模幂、素性检验、因数分解和公钥密码为例，展示数学正确性与计算效率之间的联系；密码部分只讨论可逐步核对的代数模型。

\section{输入规模、运算次数与算法描述}

正整数 $n$ 的二进制位数为 $L=\lfloor\log_2n\rfloor+1$，这是描述输入规模的自然参数。对最终非负的函数 $f,g$，且 $g$ 最终为正，记 $f(t)=O(g(t))$，表示存在常数 $C,t_0>0$ 使 $t\ge t_0$ 时 $f(t)\le Cg(t)$。常数不能随输入增长而改变。

试除到 $\sqrt n$ 的次数关于 $n$ 是平方根量级，却关于位数 $L$ 为指数级。欧几里得算法的非零余数每两步至少缩小一半：若本次余数超过前一除数的一半，下一次余数就小于该一半。因此对初始 $A\ge B\ge1$，除法次数为 $O(1+\log B)$。

快速模幂使用 $O(\log(e+1))$ 次模乘（$e\ge1$）。但一次大整数运算不能视为一个位操作。若模数有 $L$ 位、底数已约化，采用学校式乘法和长除法，则每次模乘可用 $O(L^2)$ 位操作实现，所以得到上界 $O(L^2\log(e+1))$。这是指定运算方法下的上界，并非最优复杂度断言。

全书基础算法可以依赖调用：扩展欧几里得算法给出最大公因数与贝祖系数；互素时系数给出模逆元；模逆元用于一次同余与中国剩余定理；快速模幂用于大指数运算。每步返回的是整数或剩余类，应明确其代表元范围。

例如，同为“循环 $n$ 次”，若 $n$ 是输入整数的数值，循环次数可能随输入位数指数增长；若 $n$ 是输入的位数，则含义完全不同。还应区分算法给出的有效输出与失败状态：模逆元算法在最大公因数大于 $1$ 时应报告不存在，不能把某个贝祖系数误当作逆元；同余合并算法应区分没有解与返回代表元 $0$，因为 $0$ 可能是一个完全合法的答案。

\section{素性检验与伪素数}

对奇数 $n>2$，取 $1<a<n$。若 $\gcd(a,n)>1$，已找到非平凡因子；若 $\gcd(a,n)=1$ 而 $a^{n-1}\not\equiv1\pmod n$，由费马小定理可判定 $n$ 合数。若同余成立，只能说通过了该底数的检验。通过检验的合数称为底数 $a$ 的\emph{费马伪素数（Fermat pseudoprime）}。

例如 $561=3\cdot11\cdot17$，对每个与 $561$ 互素的整数 $a$，费马小定理分别给出模 $3,11,17$ 的幂周期整除 $2,10,16$，它们均整除 $560$，故 $a^{560}\equiv1\pmod{561}$。所以仅增加若干费马检验底数不能保证正确判定素数。

\emph{Miller--Rabin 检验（Miller--Rabin primality test）}进一步使用平方根信息。对待测整数 $n$，先排除 $n<2$，对 $2,3$ 直接判为素数，对其他偶数判为合数。余下 $n\ge5$ 为奇数，写 $n-1=2^sd$，其中 $d$ 为奇数，取 $1<a<n-1$ 并先检验最大公因数。若互素，则计算 $b_0=a^d\pmod n$，随后令 $b_{j+1}=b_j^2\pmod n$。当 $b_0=1$，或某个 $0\le j<s$ 有 $b_j=-1\pmod n$ 时，该底数通过；其余情形判为合数。

该规则不会误判素数。若 $n$ 为素数，则 $b_s=1$；若 $b_0\ne1$，取第一次出现 $1$ 的位置，前一项是 $1$ 的平方根且不是 $1$，模素数时只能为 $-1$。合数则可能通过某些底数；重复抽样的概率分析需要额外定理，本节不以未经证明的概率界代替说明。

\begin{example}[检验失败给出了怎样的证据]
用底数 $2$ 对 $561$ 作 Miller--Rabin 检验。
\end{example}
\begin{solve}
$560=2^4\cdot35$，逐次平方得到
\[
 2^{35}\equiv263,\quad263^2\equiv166,\quad
 166^2\equiv67,\quad67^2\equiv1\pmod{561}.
\]
起点不是 $1$，而出现 $1$ 前没有出现 $560=-1$，所以判为合数。更具体地，$67$ 是模 $561$ 的一个不等于 $\pm1$ 的平方根 $1$；由 $561\mid(67-1)(67+1)$，求得 $\gcd(66,561)=33$、$\gcd(68,561)=17$，还找到了非平凡因子。本例说明额外检查平方根信息，能揭示普通费马检验遗漏的结构。
\end{solve}

\section{因数分解与算法问题的区别}

素性检验只需回答是否为素数，因数分解则要求给出非平凡因子或完整分解。两者的输出不同，不能把某个素性检验直接当作分解算法。

对奇合数 $n$，\emph{平方差分解（Fermat factorization）}从 $A=\lceil\sqrt n\rceil$ 开始，检查 $A^2-n$ 是否为整数平方 $B^2$；若是，则 $n=(A-B)(A+B)$，否则令 $A\leftarrow A+1$。其中 $\lceil x\rceil$ 为不小于 $x$ 的最小整数。若因子 $u,v$ 相近，则 $A=(u+v)/2$ 接近起点，因此这种方法可能较快；因子相距很远时可能很慢。例如 $6059=78^2-5^2=73\cdot83$。

进一步的 Pollard $\rho$ 方法可按如下方式试探因子：固定 $f(t)=t^2+c\pmod n$，初始化 $x=y=x_0$，反复令 $x\leftarrow f(x)$、$y\leftarrow f(f(y))$，并计算 $d=\gcd(|x-y|,n)$。若 $1<d<n$，输出因子；若 $d=n$，更换参数重试；若 $d=1$，继续迭代。实际应设置步数限额，到限额后报告此次尝试未找到因子。每个返回因子都可直接验证，但这一描述不保证任意固定参数下必定成功。

\begin{example}[观察一次因数分解迭代]
对 $n=91$ 取 $f(t)=t^2+1\pmod{91}$、$x_0=2$，执行上述迭代。
\end{example}
\begin{solve}
初始化 $x=y=2$。第一步得到 $x=5,y=26$，故 $\gcd(|5-26|,91)=\gcd(21,91)=7$，于是 $91=7\cdot13$。两个状态模 $91$ 尚不相等，但模因子 $7$ 已经相等；最大公因数把这种较小模数下的碰撞提取为因子。这个特例很快成功，并不能据此推断所有输入都只需很少几步。
\end{solve}

\section{RSA 的数论模型与正确性}

\emph{公钥密码（public-key cryptography）}区分公开参数与秘密参数。RSA 的基本数论模型如下：取不同奇素数 $p,q$，令 $n=pq$，选整数 $1<e<\varphi(n)$ 且 $\gcd(e,\varphi(n))=1$，再由扩展欧几里得算法求 $1\le d<\varphi(n)$ 使
\[
 ed\equiv1\pmod{\varphi(n)},\qquad \varphi(n)=(p-1)(q-1).
\]
对消息代表元 $0\le M<n$，密文 $C$ 为 $M^e$ 模 $n$ 的最小非负余数；解密输出 $C^d$ 模 $n$ 的最小非负余数。$(n,e)$ 为公开参数，$d$ 为私钥指数。

\begin{proposition}[RSA 的正确性]
\label{prop:rsa-correctness}
在上述条件下，对每个 $0\le M<n$ 均有 $M^{ed}\equiv M\pmod n$。
\end{proposition}
\begin{proof}
写 $ed=1+k(p-1)(q-1)$。模 $p$ 考察：若 $p\mid M$，两侧均为零；若 $p\nmid M$，由费马小定理有 $M^{p-1}\equiv1$，故 $M^{ed}\equiv M$。模 $q$ 同理，再由中国剩余定理合并。这一证明也覆盖与 $n$ 不互素的消息。
\end{proof}

\begin{example}
取 $p=5,q=11,e=3$，演示加密与解密。
\end{example}
\begin{solve}
$n=55,\varphi(n)=40$，由 $3\cdot27=81\equiv1\pmod{40}$，得 $d=27$。消息 $M=7$ 加密为 $C=7^3\bmod55=13$，解密 $13^{27}\bmod55=7$。若 $M=5$，虽不与 $55$ 互素，仍有 $C=15$ 和 $15^{27}\bmod55=5$，与定理一致。
\end{solve}

这里讨论的是基本代数模型，正确性不等于安全性。知道 $p,q$ 就能计算 $\varphi(n)$ 与私钥指数，这是一个明确的推导方向；不能据此未经证明地断言所有恢复消息的方法都等价于分解 $n$。实际密码方案还包含本章未讨论的消息编码与随机化机制。

\section{离散对数与综合实验}

在素数模数下固定原根 $g$，计算 $g^a\pmod p$ 可用快速模幂，而由结果求 $a\pmod{p-1}$ 是离散对数问题。Diffie--Hellman 的代数关系是：分别公开 $A=g^a\pmod p$、$B=g^b\pmod p$ 后，两方可分别计算
\[
 B^a\equiv A^b\equiv g^{ab}\pmod p.
\]
这证明两方得到相同的值，但没有证明第三方无法计算该值。讨论困难性时还需指定参数分布与计算模型，不能只根据运算方向不同作出结论。

\begin{example}[核对离散对数模型中的共同值]
取 $p=23,g=5$，两方分别选择指数 $a=6,b=15$，计算公开值与共同值。
\end{example}
\begin{solve}
$5^{11}\equiv-1$、$5^2\equiv2\pmod{23}$，所以 $5$ 是原根。公开值为 $A=5^6\bmod23=8$、$B=5^{15}\bmod23=19$。共同值为 $B^a\equiv A^b\equiv5^{90}$；由 $90\equiv2\pmod{22}$，它等于 $5^2\bmod23=2$。这些小参数便于演示运算，第三方也能直接列举指数找到答案；例题仅验证代数关系。
\end{solve}

\section{习题}

算法输出的含义要与定理的逻辑方向一致：发现费马或 Miller--Rabin 证据可以判合数，通过有限底数检验却不自动构成素性证明。RSA 与离散对数模型的计算验证说明正确性，效率和困难性仍需要相应的输入规模与计算模型。

\begin{enumerate}
 \item 写出快速模幂的逐步运算表，核查 $13^{27}\pmod{55}$。
 \item 验证 $561$ 对每个互素底数通过费马检验；找一个底数使它未通过 Miller--Rabin 检验。
 \item 用平方差法分解 $10403$，解释它为何适合此例。
 \item 实现扩展欧几里得、模逆元与中国剩余定理；测试无逆元和约束不相容的输入，并明确返回结果。
 \item\label{ex:algo-inverse-fail} 用贝祖等式解释为什么 $6$ 模 $15$ 没有逆元，并设计模逆元算法在此输入上的返回说明。
 \item\label{ex:algo-rsa-small} 对 $p=3,q=11,e=3$ 求 RSA 私钥指数 $d$，加密消息 $M=6$ 并解密，核对非互素消息也满足正确性。
 \item\label{ex:algo-sqrt-factor} 若奇合数 $n$ 满足 $r^2\equiv1\pmod n$ 且 $r\not\equiv\pm1\pmod n$，证明 $\gcd(r-1,n)$ 和 $\gcd(r+1,n)$ 都是非平凡因子。
 \item\label{ex:algo-test-cases} 为快速模幂写出覆盖负底数、零指数、非互素底数和非法模数的测试输入，说明各输入的预期结果或拒绝理由；再解释这些测试为何不能代替一般正确性证明。
\end{enumerate}

